\section{Introduction}


The development of safety-critical cyber-physical systems (CPS) follows an iterative process in which requirements, models, and implementations co-evolve across design cycles, driven by changing regulations, design improvements, and emerging optimization opportunities. Model-based systems engineering (MBSE)~\cite{INCOSE2023} is the established discipline for managing this complexity in safety-critical industries such as automotive and aerospace \cite{kugele2021automotive,broy2012spes,eurocae2011ed12c,rtca2011do178c}, providing a rigorous framework of formal models that capture system structure, behavior, and requirements, and through which correctness can be established and preserved across iterations.%~\cite{kugele2021automotive, broy2012spes, eurocae2011ed12c, rtca2011do178c}.

%The development of cyber-physical systems (CPS) follows an iterative process in which requirements, models, and implementations co-evolve across design cycles. Design improvements, operational adaptation and emerging optimization opportunities make this evolution necessary to ensure that the system remains fit for purpose throughout its lifetime. 

%In the context of safety-critical and real-time CPS, this iterative co-evolution required a rigorous engineering discipline to ensure that the system's design and implementation remain correct with respect to its requirements across iterations. To this end, model-based systems engineering (MBSE) \cite{INCOSE2023} is a well-established mechanism for developing and evolving such systems \cite{kugele2021automotive,broy2012spes,eurocae2011ed12c,rtca2011do178c}, providing a well-understood framework centered on formal models that capture the system's structure, behavior, and requirements.

Complementary to MBSE, the concept of Digital Twins (DTs) \cite{grieves2005} has emerged as a powerful software engineering paradigm to support decision-making across the lifecycle of CPSs. A DT pairs a virtual representation of the system with a connection to the deployed system, referred to as the Physical Twin (PT), enabling continuous synchronization, run-time monitoring, and decision support throughout the system’s operational life~\cite{grieves2005,Tao2019, Minerva2020}. DTs have demonstrated broad applicability across domains~\cite{Woodcock2021, Liu2021}, and their use has been extended beyond monitoring and predictive analytics to support iterative design refinement, virtual prototyping, and post-deployment adaptation~\cite{Tao2018, Mertens2024, Lehner2025}.


%These properties hinge on the requirement that the DT and PT remain semantically aligned across operation and evolution, such that the DT remains a faithful and coherent representation of the PT’s structure and behavior. 

These properties hinge on the requirement that he meaning assigned to observed data must agree between the two twins, and their observable behaviors must correspond, i.e. that the DT and PT remain semantically aligned across operation and evolution. Absence or loss of either form of alignment undermines the correctness of any DT-based reasoning, as a DT that misinterprets a sensor reading or exhibits behavior inconsistent with the PT cannot be trusted for monitoring or decision-making. This is particularly critical in safety-critical domains, and is further exacerbated in modern federated architectures such as Gaia-X~\cite{Gaiax2020}, where independently developed DTs must interoperate and reason jointly.


%that the DT and PT remain semantically aligned across operation and evolution: the meaning assigned to observed data must agree between the two twins, and their observable behaviors must correspond. Absence or loss of either form of alignment undermines the correctness of any DT-based reasoning — a DT that misinterprets a sensor reading or exhibits behavior inconsistent with the PT cannot be trusted for monitoring or decision-making. This is particularly critical in safety-critical domains, and is further exacerbated in modern federated architectures such as Gaia-X~\cite{GaiaX2020}, where independently developed DTs must interoperate and reason jointly.


%These properties hinge on the requirement that the DT remains a faithful and coherent representation of the PT’s structure and behavior, i.e., that the they remain semantically aligned across operation and evolution. Absence or loss of alignment undermines the correctness of any DT-based reasoning, which is particularly critical in safety-critical domains, and it is further exacerbated in modern federated architectures such as Gaia-X~\cite{GaiaX2020}, where independently developed DTs must interoperate and reason jointly.

Current approaches to DT-PT alignment, however, address neither requirement fully. Runtime monitors~\cite{10.1145/3652620.3688267} are bounded by finite observations and cannot certify the unobserved behavioral space. Model synchronisation mechanisms~\cite{lehner2021aml4dt, Mertens2024} ensure structural correspondence but not semantic equivalence. Consistency frameworks~\cite{kamburjan2022twinningbyconstruction, Gil2025} check only an explicitly chosen property fragment, leaving the remainder unchecked. Ontology-based matching~\cite{10.1016/j.rcim.2023.102649, KARABULUT2024442} aligns static data schemas rather than dynamic behavior, meaning DT and PT can conform to the same schema while exhibiting observably different behaviors under it. No existing mechanism therefore establishes, at design time, that the DT and PT agree on the meaning of observed data and on their observable behavior, nor guarantee that this agreement is preserved as both evolve across design iterations and operational use.

%Current approaches to pursue DT-PT alignment, however, do not provide such guarantees.
%Runtime monitors~\cite{10.1145/3652620.3688267} are bounded by finite observations and cannot certify the unobserved behavioral space. Model synchronisation mechanisms~\cite{lehner2021aml4dt, Mertens2024} ensure structural correspondence but not semantic equivalence. Consistency frameworks~\cite{kamburjan2022twinningbyconstruction, Gil2025} check only an explicitly chosen property fragment, leaving the remainder unchecked. Ontology-based matching~\cite{10.1016/j.rcim.2023.102649, KARABULUT2024442} aligns static data schemas rather than dynamic behavior, meaning DT and PT can conform to the same schema while exhibiting observably different behaviors under it. No existing mechanism therefore establishes at design time that the DT and PT are semantically aligned both at the data and behavior level, nor preserve this alignment as both evolve across design iterations and operational use.
%
%%that this alignment is preserved as both evolve across design iterations and operational use.
%

%No mechanism therefore establishes, at design time, that the DT and PT are behaviorally equivalent with respect to a shared domain semantics, nor that this equivalence is preserved as both evolve.


To bridge this gap, we present a formal software engineering framework based on MBSE for semantically aligned DTs. We formally characterize semantic alignment via a shared ontology between PT and DT, establish it once and compositionally at design time, and derive conditions under which it is preserved across iterative co-evolution. We provide an algorithm to check alignment, and prove that runtime monitors and ontology-based matchers are fundamentally incomplete approximations of this condition.  The framework delivers these guarantees at no additional engineering cost: the ontology is already required for Gaia-X compliance, and a single design-time verification replaces the indefinite runtime overhead of monitor instrumentation, deployment, and maintenance.

%To bridge this gap, we present a formal software engineering framework for semantically aligned DTs under MBSE. We first provide a formal characterization of semantic alignment based on a shared ontology between PT and DT. This is established once and compositionally at design time, and establish the condition under which it is preserved across iterative co-evolution. We provide an algorithm to establish and maintain this alignment and show that monitors and ontology-based matchers are fundamentally incomplete approximations of this condition. 

%We define semantic alignment as weak timed bisimulation over a shared ontological action alphabet, establish it at design time once and compositionally, and prove it is preserved at runtime and across iterative co-evolution. We further show that existing runtime approaches are fundamentally incomplete approximations of this condition, as they operate over finite observations and selected properties, while bisimulation characterizes full behavioral equivalence.


%To bridge this gap, we present a formal software engineering framework for developing semantically aligned DTs under MBSE. Grounded in the theory of refinement and bisimulation, we define semantic alignment as weak timed bisimulation over a shared ontological action alphabet established by a domain ontology, which we generalize as first-order logic (FOL) to encompass existing ontology languages and knowledge graph approaches. We then leverage this concept to establish alignment at design time completely and compositionally, which is guaranteed to be preserved at runtime and provides a mechanism to establish and re-establish alignment at design time, across iterative development and evolution. We also provide conditions under which the ontology can be evolved without invalidating previously established alignment results, and we show that existing runtime approaches are fundamentally incomplete approximations of this condition, as monitors and ontology-based matchers operate over finite observations and selected properties, while bisimulation characterizes full behavioral equivalence. 


%grounded in first-order logic, which generalizes existing ontologies concepts based on knowledge graphs. 


%as well as across iterative development and evolution. We further prove that existing runtime approaches are fundamentally incomplete approximations of this condition.

%Through this concept, our framework establishes alignment at design time once, completely, and compositionally, 


%To bridge this gap, we present a formal framework for semantic alignment of Digital Twins under MBSE, grounded in weak timed bisimulation over a shared ontological action alphabet. Our framework establishes alignment at design time — once, completely, and compositionally — by relating PT and DT implementations to their respective timed automata views via simulation refinement, a condition guaranteed by MBSE development discipline. We further prove that existing runtime approaches are fundamentally incomplete approximations of this condition: monitors and ontology-based matchers operate over finite observations and selected properties, while bisimulation characterizes full behavioral equivalence via Hennessy-Milner logic. We evaluate the framework on a smart farming case study, comparing against five representative approaches and demonstrating scenarios where existing approaches fail and ours succeeds.


%Ontologies are used to represent domain knowledge and establish a shared semantic interface between the DT and PT in the form of weak

%In this paper, we present a formal software engineering framework for developing semantically aligned DTs under MBSE. Grounded in the theory of refinement and bisimulation, our approach provides a rigorous methodology for establishing and preserving semantic alignment between the DT and PT across design iterations, operational use, and system evolution. Following the literature, we use ontologies (which we generalize as FOL) to establish a form of weak bisimualtion.... 

%Our approach is grounded in the theory of refinement, which we 


%Yet realizing this potential requires that the DT remains semantically aligned with the PT across development, operation, and evolution — that is, that the DT's models faithfully reflect the structure and behavior of the physical system it represents. In safety-critical domains, misalignment is not merely an engineering inconvenience: it directly undermines the reliability of monitoring, diagnosis, and autonomous decision-making. Beyond individual systems, semantic alignment is equally a prerequisite for interoperability in federated architectures~\cite{GaiaX2020}, where DTs developed independently by different parties must compose and reason jointly — a property that run-time monitors alone cannot guarantee. Despite its importance, DT development under MBSE currently lacks a formal mechanism for establishing and preserving semantic alignment across the coupled evolution of the DT and its PT.

%In order to realize the full potential of DTs, however, a critical challenge remains: ensuring that the DT and PT remain semantically aligned across development, operation, and evolution. Semantic alignment means that the DT's model and data accurately reflect the PT's behavior and state, enabling reliable monitoring, analysis, and decision-making. This is particularly crucial in safety-critical domains, where misalignment can lead to incorrect diagnoses, unsafe decisions, and ultimately harm to people or the environment. Moreover, in federated architectures like Gaia-X~\cite{GaiaX2020}, where multiple DTs developed by different parties must interoperate, semantic alignment becomes a prerequisite for interoperability and trustworthiness.
%DT development under MBSE, as currently practiced, lacks a formal mechanism for preserving this correctness. A DT is not a further-refined implementation of the PT but a parallel artifact with its own purpose (monitoring and decision support) which co-exist and co-evolve with the PT. This duality, combined with iterative development, creates a coupled correctness problem that manifests across multiple dimensions.

%DT development under MBSE, as currently practiced, lacks a formal mechanism for preserving correctness across iterative design and evolution. For traditional MBSE, this mechanism has been the concept of refinement, which captures when one component is a more detailed realization of another, thereby yielding desirable properties such as correctness preservation, compositionality, and substitutability. Unlike traditional MBSE, however, where refinement relates successive development stages of a single artifact, DT-based workflows introduce a dual structure, where a DT is not a further-refined implementation of the PT but a parallel artifact with its own purpose (monitoring and decision support) which co-exist and co-evolve with the PT. This duality, combined with iterative development, creates a coupled correctness problem that manifests across multiple dimensions.



%However, introducing refinement into the twin setting raises several design choices that standard refinement-based development does not encounter. 




%At the specification level, the DT and PT serve different purposes(the PT capturing required system behavior, and the DT supporting monitoring and prediction) but they cannot be disjoint: the DT must consistently reflect the PT on a shared subset of events. At design time, semantic alignment between DT and PT must be established, as any drift undermines the trustworthiness of any DT-based monitoring and decision-making. Finally, during system evolution and design iterations, updates proposed via the DT must preserve specification-level correctness, maintain alignment with the PT, and remain deployable under platform constraints.